% Autor: Elias Welt
\section{Detaillierte Analyse ausgewählter Module}
Um die technische Tiefe der Implementierung zu verdeutlichen, werden im Folgenden drei Kernmodule detailliert anhand ihres Quellcodes analysiert.

\subsection{Quiz-Editor Logik (QuizModal.jsx)}
Der Quiz-Editor ermöglicht das Erstellen und Bearbeiten von verschachtelten Datenstrukturen (Quiz -> Fragen -> Antworten). Die Herausforderung hierbei ist das State-Management einer mehrdimensionalen Struktur.

\subsubsection{Zustandsverwaltung}
Der State \texttt{questions} ist ein Array von Objekten. Jedes Objekt repräsentiert eine Frage und enthält wiederum ein Array von \texttt{answers}.

\begin{lstlisting}[language=JavaScript, caption={Initialer State im QuizModal}]
const [questions, setQuestions] = useState([
  { 
    text: '', 
    answers: [{ text: '', isCorrect: false }] 
  }
]);
\end{lstlisting}

\subsubsection{Dynamische Formular-Manipulation}
Beim Ändern eines Antworttextes darf der State nicht direkt angepasst werden. Stattdessen muss eine Deep Copy der betroffenen Ebenen erstellt werden.
\texttt{handleAnswerChange} demonstriert dabei die Immutability~\cite{ReactImmutability} von React:

\begin{lstlisting}[language=JavaScript, caption={Update-Logik für verschachtelte Antworten}]
const handleAnswerChange = (questionIndex, answerIndex, event) => {
    const newQuestions = [...questions];
    newQuestions[questionIndex].answers[answerIndex].text = event.target.value;
    setQuestions(newQuestions);
};
\end{lstlisting}

\subsubsection{Validierung und Speichern}
Beim Speichern (\texttt{handleSubmit}) werden die Daten an die API gesendet. Hierbei wird unterschieden, ob ein neues Quiz erstellt (\texttt{POST}) oder ein bestehendes bearbeitet (\texttt{PUT}) wird.

\subsection{Flashcards: 3D-Interaktion und State}
Die Komponente \texttt{Flashcards.jsx} verwendet eine interessante Datenstruktur für den Zustand der umgedrehten Karten. Anstatt für jede Karte einen Boolean im Objekt zu speichern (was eine Änderung des Datenmodells erfordern würde), wird ein \texttt{Set} verwendet, das die IDs der umgedrehten Karten speichert.

\begin{lstlisting}[language=JavaScript, caption={Nutzung eines Sets für UI-State}]
const [flippedCards, setFlippedCards] = useState(new Set());

const flipCard = (id) => {
    setFlippedCards(prevFlippedCards => {
        const newFlippedCards = new Set(prevFlippedCards);
        if (newFlippedCards.has(id)) {
            newFlippedCards.delete(id);
        } else {
            newFlippedCards.add(id);
        }
        return newFlippedCards;
    });
};
\end{lstlisting}

Diese Methode ist performant und trennt den UI-Zustand (ist die Karte gedreht?) sauber von den eigentlichen Daten (Frage/Antwort). Beim Rendern wird \texttt{flippedCards.has(card.id)} geprüft, um die CSS-Klasse \texttt{.flipped} bedingt hinzuzufügen.
\newpage

\subsection{Admin Panel: Kaskadierende Datenabfragen}
Im \texttt{AdminPanel.jsx} muss beim Klick auf einen Benutzer dessen Dateiliste geladen werden. Dies ist ein klassisches "Master-Detail"-Szenario.
Um unnötigen Netzwerktraffic zu vermeiden, werden die Dateien nicht initial für alle Benutzer geladen ("Lazy Loading").

\begin{lstlisting}[language=JavaScript, caption={Asynchrones Laden von Detaildaten}]
const handleUserClick = (user) => {
    setSelectedUser(user);
    setError('');
    setEditUserMode(false);
    fetchUserFiles(user.userId);
};

const fetchUserFiles = async (userId) => {
    try {
        setLoading(true);
        const files = await api.get(`/Users/${userId}/documents`);
        setUserFiles(files);
    } catch (err) {
        setError('Could not load user files.');
    } finally {
        setLoading(false);
    }
};
\end{lstlisting}

Diese Implementierung sorgt für eine schnelle initiale Ladezeit der Benutzertabelle, da die (potenziell große) Menge an Dateimetadaten erst bei Bedarf ("`On-Demand"') abgerufen wird.
\newpage
